%% GRPO Importance Ratio 深度排查笔记
%% 记录时间：2025-04-11
%% 背景：R3（BAGEL 图像编辑 RL 训练）中关于 GRPO importance ratio 的实现方式引发的疑问

\section{GRPO Importance Ratio 深度排查}

\note{记录时间：2025-04-11 \quad 背景：R3（BAGEL 图像编辑 RL 训练）中关于 GRPO importance ratio
的实现方式引发的疑问，经过查阅源码、GitHub 讨论和论文，最终得出结论。}

% ============================================================
\subsection{起点：疑问}
% ============================================================

在 \code{train/online\_rl.py} 中，image policy loss 的写法是：

\begin{lstlisting}[language=Python]
img_per_token_loss = -(
    torch.exp(log_prob_per_image - log_prob_per_image.detach()) * cur_image_advantages
).sum() / (1e-4 + total_image_num)
\end{lstlisting}

文本 policy loss 也类似：

\begin{lstlisting}[language=Python]
per_token_loss = -(
    torch.exp(per_token_logps - per_token_logps.detach()) * cur_text_advantages
).sum() / (1e-4 + total_text_grad_tokens)
\end{lstlisting}

\begin{warnbox}
\textbf{疑问}：\code{log\_prob - log\_prob.detach()} 永远等于 $0$，$\exp(0) = 1$，
那 ratio 不就永远是 $1$ 吗？\code{.detach()} 怎么能作为 old policy？
\end{warnbox}

% ============================================================
\subsection{对比参考：flow\_grpo 的实现}
% ============================================================

\note{来源：\texttt{https://github.com/yifan123/flow\_grpo}}

\code{flow\_grpo} 的做法分为两个阶段：

\begin{lstlisting}[language=Python]
# 阶段 1：rollout 时保存 old log_prob（权重为 theta_0）
images, latents, log_probs = pipeline_with_logprob(...)
samples["log_probs"] = log_probs   # <- 存入内存，这才是真正的 old policy

# 阶段 2：多次梯度更新（theta 从 theta_0 变成 theta_1, theta_2, ...）
for epoch in range(num_epochs):
    log_prob_new = recompute_log_prob(latents)   # 当前权重重新算
    ratio = torch.exp(log_prob_new - sample["log_probs"])  # pi_new / pi_old，真的不等于 1

    unclipped = -advantages * ratio
    clipped   = -advantages * torch.clamp(ratio, 1-eps, 1+eps)
    loss = torch.mean(torch.maximum(unclipped, clipped))
\end{lstlisting}

\begin{conceptbox}
\textbf{\code{flow\_grpo} 的关键区别}：old policy 是 rollout 时真正存下来的 $\theta_0$
的 log\_prob。多次梯度更新后 $\theta$ 发生漂移，ratio 真实地偏离 $1$，clip 具有实际意义。
\end{conceptbox}

% ============================================================
\subsection{社区讨论的结论}
% ============================================================

\subsubsection{来源 1：TRL PR \#2565}

\note{\texttt{https://github.com/huggingface/trl/pull/2565\#issuecomment-2595837761}}

TRL 作者 \code{qgallouedec} 的回答：

\begin{intubox}
\textit{``The math is correct; if you look at the loss, it is indeed equal to the KL value.
However, in terms of differentiation, we cannot remove the ratio term justifying that it
equals 1.''}

他指出，必须保留 ratio 项的形式：
\[
    \frac{\pi_\theta}{\lfloor \pi_\theta \rfloor_{\mathrm{sg}}}
\]
其中下标 $\mathrm{sg}$ 表示 stop-gradient（即 \code{.detach()}）。

写成 \code{exp(logp - logp.detach())} 的原因：\textbf{不是为了模拟 old policy，
而是为了保留梯度通道}。
\end{intubox}

\begin{warnbox}
如果直接删掉 ratio，loss 变成常数（$-\text{advantage}$），梯度为零，\warn{模型不会更新}。
\end{warnbox}

\subsubsection{来源 2：R1-V Issue \#174}

\note{\texttt{https://github.com/StarsfieldAI/R1-V/issues/174}}

有人认为这是 bug（``open-r1 库一开始犯了这个错误，现在已纠正''），有人认为不是。

\textbf{结论}：两种实现都是正确的，只是适用场景不同：

\begin{center}
\begin{tabular}{>{\bfseries}p{4.2cm} p{4.8cm} p{4.8cm}}
\hline
\rule{0pt}{2.4ex}
& \textbf{\code{logp.detach()} 方式（TRL/R3）}
& \textbf{存储真实 old logp 方式（flow\_grpo / open-r1 修正版）} \\
\hline
每次 rollout 做几次梯度更新 & \good{1 次} & \good{多次}（\code{num\_epochs}）\\
old policy 来源 & 当前 forward 的 stop-grad 拷贝 & rollout 时真正存下来的 $\theta_0$ \\
ratio 数值 & 永远 $= 1$ & 真的偏离 $1$ \\
clip 有用吗 & \bad{无用}（形同虚设） & \good{有用} \\
数学自洽吗 & \good{✓} 单步更新下自洽 & \good{✓} 多步更新下自洽 \\
\hline
\end{tabular}
\end{center}

\subsubsection{来源 3：TRL Issue \#2608（最终结论）}

\note{\texttt{https://github.com/huggingface/trl/issues/2608}}

\code{qgallouedec} 明确说明：

\begin{intubox}
\textit{``In the current implementation, we're just update once after a generation.
In fact, we align with this sentence from the paper: \emph{The policy model only has a
single update following each exploration stage.} Therefore it implies that
$\pi_{\theta_{\mathrm{old}}} = \pi_\theta$, and the equation can be simplified to\ldots''}
\end{intubox}

\[
    \mathcal{J}_{\mathrm{GRPO}}(\theta)
    = \frac{1}{G} \sum_{i=1}^{G} \frac{1}{|o_i|} \sum_{t=1}^{|o_i|}
      \left[
          \frac{\pi_\theta}{\lfloor \pi_\theta \rfloor_{\mathrm{sg}}} \,\hat{A}_{i,t}
          \;-\;
          \beta\, \mathbb{D}_{\mathrm{KL}}\!\left[\pi_\theta \,\|\, \pi_{\mathrm{ref}}\right]
      \right]
\]

用户 \code{ghrua} 进一步指出（\key{关键结论}）：

\begin{summarybox}
\textit{``I think the two lines of code are equivalent\ldots the single-update GRPO
degrades to the standard policy gradient, which is consistent with Eq.~(17) in
DeepSeekMath.''}

\medskip
\begin{lstlisting}[language=Python]
# 这两行完全等价（梯度相同）
loss = -torch.exp(per_token_logps - per_token_logps.detach()) * advantages
loss = -per_token_logps * advantages   # <- 就是 REINFORCE
\end{lstlisting}

\medskip
\textbf{证明（链式法则）：}
\[
    \frac{\partial}{\partial\theta}
    \Bigl[\exp\!\bigl(\log p - \underbrace{\log p}_{\mathrm{sg}}\bigr)\Bigr]
    = \exp(0) \cdot \frac{\partial \log p}{\partial \theta}
    = \frac{\partial \log p}{\partial \theta}
\]
\end{summarybox}

% ============================================================
\subsection{最终结论}
% ============================================================

\subsubsection{R3 的 GRPO loss 本质上是什么？}

\begin{summarybox}
\textbf{带 group-relative advantage 归一化的 REINFORCE（Policy Gradient）}
\[
    \mathcal{L} = -\log p_\theta \cdot A
\]
其中 advantage $A$ 用 GRPO 的 group-relative 方式归一化（组内减均值除标准差）。
\end{summarybox}

\subsubsection{三个层次的理解}

\begin{center}
\begin{tabular}{>{\bfseries}p{2.5cm} p{10cm}}
\hline
\rule{0pt}{2.4ex}
层次 & 结论 \\
\hline
数值上 & \code{exp(logp - logp.detach())} $= 1$，ratio 永远等于 $1$ \\
梯度上 & 等价于直接写 $\log p \times A$，即标准 REINFORCE \\
算法上 & 单步 GRPO $=$ REINFORCE $+$ GRPO 的 group-relative advantage \\
\hline
\end{tabular}
\end{center}

\subsubsection{R3 和 flow\_grpo 的本质区别}

\begin{center}
\begin{tabular}{>{\bfseries}p{3.8cm} p{4.8cm} p{4.8cm}}
\hline
\rule{0pt}{2.4ex}
& \textbf{R3（单步）} & \textbf{flow\_grpo（多步）} \\
\hline
算法名称 & REINFORCE + group advantage & GRPO with importance sampling \\
采样效率 & \bad{低}（每个 rollout 只用 1 次） & \good{高}（每个 rollout 复用 \code{num\_epochs} 次）\\
实现复杂度 & 简单 & 需要存储 \code{old\_log\_prob}，处理多轮更新 \\
是否需要 clip & \bad{不需要}（ratio$=1$，clip 无效） & \good{需要}（ratio 真的偏离 $1$）\\
是否正确 & \good{✓} 单步下数学自洽 & \good{✓} 多步下数学自洽 \\
\hline
\end{tabular}
\end{center}

\subsubsection{ref\_model 的角色（常见混淆点）}

\begin{warnbox}
\code{ref\_model} \warn{不是} old policy，和 importance ratio 没有关系。

\begin{itemize}
    \item \key{old policy}：用于 importance ratio（$\pi_{\mathrm{new}} / \pi_{\mathrm{old}}$），
          在单步实现中退化为 \code{logp.detach()}
    \item \key{ref\_model}：用于 KL penalty（防止训练后的策略偏离预训练模型太远），
          是一个\textbf{冻结的预训练权重快照}
\end{itemize}
\end{warnbox}

\begin{lstlisting}[language=Python]
# KL penalty（ref_model 的用途）
kl = exp(ref_logp - logp) - (ref_logp - logp) - 1   # K3 估计

# importance ratio（和 ref_model 无关）
ratio = exp(logp - logp.detach())   # = 1
\end{lstlisting}

% ============================================================
\subsection{如果想升级为真正的多步 GRPO}
% ============================================================

\begin{conceptbox}
R3 中 \code{RolloutStepResult} 已经存了 \code{log\_prob}，基础设施完备。
以下改动是\textbf{性能优化方向}（提高采样效率），而不是纠正错误。
\end{conceptbox}

\textbf{步骤 1：rollout 时标记 old\_log\_prob}（已有，无需改动）

\begin{lstlisting}[language=Python]
rollout_result.log_prob = log_probs[i]   # 这就是 old log_prob
\end{lstlisting}

\textbf{步骤 2：训练循环中重新计算 new log\_prob}

\begin{lstlisting}[language=Python]
for epoch in range(num_epochs):   # 新增多次更新
    log_prob_new = recompute_from_packed_input(...)
    old_log_prob = loss_info["old_log_probs"]   # 从 rollout 存的值取

    ratio = torch.exp(log_prob_new - old_log_prob)   # 真实 ratio

    unclipped = -advantages * ratio
    clipped   = -advantages * torch.clamp(ratio, 1 - epsilon, 1 + epsilon)
    loss = torch.mean(torch.maximum(unclipped, clipped))
\end{lstlisting}

% ============================================================
\subsection{补充：什么是 REINFORCE（Policy Gradient）}
% ============================================================

\subsubsection{背景}

REINFORCE 是 Williams（1992）提出的最基础的策略梯度算法，是所有现代强化学习方法
（PPO、GRPO 等）的祖先。

\subsubsection{核心思想}

智能体（策略 $\pi_\theta$）采取行动 $a$，获得奖励 $R$。目标是最大化期望奖励：
\[
    J(\theta) = \mathbb{E}_{\tau \sim \pi_\theta}\bigl[R(\tau)\bigr]
\]

对 $\theta$ 求梯度，利用 \key{log-derivative trick}：
\[
    \nabla_\theta J(\theta)
    = \mathbb{E}\!\left[R(\tau) \cdot \nabla_\theta \log \pi_\theta(a \mid s)\right]
\]

\begin{intubox}
\textbf{直觉}：如果某个行动得到了高奖励，就增大它的概率（$\log\pi$ 增大）；
如果得到了低奖励，就减小它的概率。
\end{intubox}

\subsubsection{对应的 loss 写法（PyTorch 风格）}

\begin{lstlisting}[language=Python]
loss = -(log_prob * reward).mean()
#        ^  要最大化 J，等价于最小化 -J
#        ^  log pi_theta(a|s)
#                   ^  奖励 R（或 advantage A）
\end{lstlisting}

\subsubsection{Advantage 版本（减 baseline）}

原始 REINFORCE 方差很大。用 advantage $A = R - b$（$b$ 是 baseline，通常是 value
function 的预测值）代替原始奖励，可以大幅降低方差：

\begin{lstlisting}[language=Python]
loss = -(log_prob * advantage).mean()
\end{lstlisting}

\begin{conceptbox}
\begin{itemize}
    \item $A > 0$：这个行动比平均水平\good{好} \quad$\Rightarrow$\quad 增大概率
    \item $A < 0$：这个行动比平均水平\bad{差} \quad$\Rightarrow$\quad 减小概率
    \item $A = 0$：和平均水平一样 \quad$\Rightarrow$\quad 不更新
\end{itemize}
\end{conceptbox}

\subsubsection{GRPO 的 advantage 怎么算}

GRPO 不需要 Critic（value function），用\key{组内相对奖励}作为 advantage：

\begin{lstlisting}[language=Python]
# 同一个 prompt，采样 G 个输出，计算相对排名
rewards = [r_1, r_2, r_3, ..., r_G]             # G 个输出的奖励
mean_r  = mean(rewards)
std_r   = std(rewards)
advantage_i = (r_i - mean_r) / (std_r + 1e-8)  # 组内 z-score 归一化
\end{lstlisting}

不需要额外训练一个 Critic 网络，结构简单，这是 GRPO 相比 PPO 的主要优势。

\subsubsection{REINFORCE $\to$ PPO $\to$ GRPO 演化路径}

\begin{center}
\begin{tabular}{>{\bfseries}p{3.2cm} p{5.0cm} p{5.0cm}}
\hline
\rule{0pt}{2.4ex}
算法 & 改进点 & 缺陷 \\
\hline
REINFORCE (1992)
    & 最基础的策略梯度
    & 单步更新，采样效率低；方差大 \\
PPO (2017)
    & importance ratio 让同一批数据多次更新；clip 防止更新步长过大
    & 需要 Critic（value function），参数量翻倍，训练复杂 \\
GRPO (2024, DeepSeekMath)
    & 去掉 Critic，用 group-relative reward 代替 advantage；保留 importance ratio + clip（多步更新时）
    & 单步退化：$\pi_{\mathrm{old}} = \pi_{\mathrm{new}} \Rightarrow \text{ratio} = 1
      \Rightarrow$ 等价于 REINFORCE + group advantage \\
\hline
\end{tabular}
\end{center}

\subsubsection{对应到 R3 代码}

\begin{lstlisting}[language=Python]
# R3 的实际计算（单步更新时）
loss = -(log_prob_per_image * cur_image_advantages).sum() / total_image_num

# 其中 cur_image_advantages 来自 GRPO 的 group-relative 归一化：
# advantage_i = (reward_i - mean(rewards_in_group)) / std(rewards_in_group)
\end{lstlisting}

\begin{summarybox}
这就是完整的 \key{REINFORCE with group-relative baseline}，数学上等价于单步 GRPO。
\end{summarybox}

% ============================================================
\subsection{参考链接}
% ============================================================

\begin{itemize}
    \item TRL PR \#2565 作者解释：
          \texttt{https://github.com/huggingface/trl/pull/2565\#issuecomment-2595837761}
    \item TRL Issue \#2608 最终讨论：
          \texttt{https://github.com/huggingface/trl/issues/2608}
    \item R1-V Issue \#174 社区争论：
          \texttt{https://github.com/StarsfieldAI/R1-V/issues/174}
    \item flow\_grpo 多步实现参考：
          \texttt{https://github.com/yifan123/flow\_grpo}
    \item DeepSeekMath 论文 Eq.~(17)：单步 GRPO 退化为标准 policy gradient
\end{itemize}
